\part{Geometric Mechanics and Rigid-Body Flight}
\section{Configuration, velocity, and degrees of freedom}
\begin{definition}[Rigid body; Goldstein 3e, \S\S1.2--1.3, pp.~11--12,
and \S4.1, p.~134 \cite{dgdt:GPS2001}]\label{gnc:def:rigid-body}
A rigid body is a system of $N$ mass points, with masses $m_i>0$ and
inertial positions $\mathbf r_i(t)\in\mathbb R^3$, subject for all $t$ to
the holonomic constraints
\begin{equation}\label{gnc:rigid-constraints}
 f_{ij}(\mathbf r_1,\dots,\mathbf r_N)
 =(\mathbf r_i-\mathbf r_j)\cdot(\mathbf r_i-\mathbf r_j)-c_{ij}^2=0,
 \qquad 1\le i<j\le N,
\end{equation}
with constants $c_{ij}=c_{ji}>0$. Its configuration space is
\begin{equation}\label{gnc:rigid-level-set}
 Q_{\mathrm{rigid}}
 =\bigl\{(\mathbf r_1,\dots,\mathbf r_N)\in\mathbb R^{3N}:
   \lvert\mathbf r_i-\mathbf r_j\rvert=c_{ij},\ 1\le i<j\le N\bigr\}.
\end{equation}
\end{definition}
Goldstein writes Eq.~\eqref{gnc:rigid-constraints} as
$(\mathbf r_i-\mathbf r_j)^2-c_{ij}^2=0$ (p.~12) and as Eq.~(4.1),
$r_{ij}=c_{ij}$, where $r_{ij}=\lvert\mathbf r_i-\mathbf r_j\rvert$.

A point mass in Euclidean space has configuration space $Q=\mathbb R^3$.
A rigid body with one material point fixed has configuration space $SO(3)$,
whether or not that point is the center of mass. Choose a right-handed
orthonormal body frame at that point and a right-handed inertial frame.
Our rotation $R$ maps body components to inertial components:
\begin{equation}\label{gnc:configuration-map}
 \mathbf x_I=R\mathbf x_B,\qquad R^TR=I,\qquad \det R=1.
\end{equation}
The columns of $R$ are the body basis vectors expressed in the inertial
frame; the inverse component transformation is $\mathbf x_B=R^T\mathbf x_I$.
We distinguish configurations by a fixed material frame, even if the
body's shape or inertia has a symmetry.
Section~\ref{quat:sec:orthogonal-submanifold} proves that these rotations
form a three-dimensional embedded submanifold of $\mathbb R^9$.
Its regular-value argument justifies the count $9-6=3$.
The motion $t\mapsto R(t)$ is a curve on this manifold, rather than an
arbitrary curve through nine independent matrix entries.
A freely
moving rigid body has configuration space $\mathbb R^3\times SO(3)$,
identified with $SE(3)$ as a manifold. Its group law is
\begin{equation}\label{gnc:group-law}
 (\mathbf r,R)(\mathbf s,S)=(\mathbf r+R\mathbf s,RS).
\end{equation}
The rotation acts on the translation in Eq.~\eqref{gnc:group-law}; this is
why the rigid-motion group is a semidirect product.

Rigid motions preserve the constraints of
Definition~\ref{gnc:def:rigid-body}. Fix body components $\mathbf s_i$ of
the mass points, satisfying
$\lvert\mathbf s_i-\mathbf s_j\rvert=c_{ij}$, and set
$\mathbf r_i=\mathbf r+R\mathbf s_i$. The translation cancels in the
difference, and $R^TR=I$ removes the rotation:
\begin{align*}
 (\mathbf r_i-\mathbf r_j)\cdot(\mathbf r_i-\mathbf r_j)
 &=\bigl(R(\mathbf s_i-\mathbf s_j)\bigr)^T R(\mathbf s_i-\mathbf s_j)\\
 &=(\mathbf s_i-\mathbf s_j)^TR^TR(\mathbf s_i-\mathbf s_j)\\
 &=c_{ij}^2.
\end{align*}
The calculation uses only $R^TR=I$, so an improper orthogonal matrix
($\det R=-1$) also satisfies Eq.~\eqref{gnc:rigid-constraints}; it gives the
mirror-image configuration. A continuous motion $R(t)$ starting at $R(0)=I$
has $\det R(t)$ continuous with values $\pm1$, hence $\det R(t)=1$.
This is why the configuration map~\eqref{gnc:configuration-map} imposes
$\det R=1$.

Degrees of freedom count independent configuration coordinates. The
three translational and three rotational degrees of freedom give twelve
independent position--velocity state variables for a second-order rigid-body
model. A unit quaternion uses four components to represent three rotational
degrees of freedom, subject to one scalar constraint. A translational
three-DOF flight model omits attitude dynamics under additional modeling
assumptions; it does not follow merely by invoking symmetry reduction.

\subsection{Direction cosines}\label{gnc:sec:direction-cosines}
Let $\mathbf e_1,\mathbf e_2,\mathbf e_3$ be the inertial basis and
$\mathbf b_1,\mathbf b_2,\mathbf b_3$ the body basis of
Eq.~\eqref{gnc:configuration-map}, both orthonormal and right-handed.
The direction cosines are
\begin{equation}\label{gnc:direction-cosines}
 R_{ij}=\mathbf e_i\cdot\mathbf b_j
 =\cos\angle(\mathbf e_i,\mathbf b_j),\qquad i,j=1,2,3;
\end{equation}
the second equality holds because both are unit vectors. The nine
$R_{ij}$ form a real $3\times3$ matrix, $R\in M_3(\mathbb R)$.
Expand each basis in the other, with coefficients $\beta_j^{\,\,k}$ and
$\gamma_i^{\,\,k}$, and dot with a basis vector:
\begin{align*}
 \mathbf b_j=\beta_j^{\,\,k}\mathbf e_k:&\qquad
 \mathbf e_i\cdot\mathbf b_j
 =\beta_j^{\,\,k}\,\mathbf e_i\cdot\mathbf e_k
 =\beta_j^{\,\,k}\delta_{ik}=\beta_j^{\,\,i},\\
 \mathbf e_i=\gamma_i^{\,\,k}\mathbf b_k:&\qquad
 \mathbf b_j\cdot\mathbf e_i
 =\gamma_i^{\,\,k}\,\mathbf b_j\cdot\mathbf b_k
 =\gamma_i^{\,\,k}\delta_{jk}=\gamma_i^{\,\,j}.
\end{align*}
Both left-hand sides are $R_{ij}$ by Eq.~\eqref{gnc:direction-cosines},
so $\beta_j^{\,\,i}=R_{ij}$ and $\gamma_i^{\,\,j}=R_{ij}$:
\begin{equation}\label{gnc:dcm-basis}
 \mathbf b_j=R_{ij}\mathbf e_i,\qquad
 \mathbf e_i=R_{ij}\mathbf b_j .
\end{equation}

Let a vector be $\mathbf a=a_I^i\mathbf e_i=a_B^j\mathbf b_j$;
$a_I^i$ is component $i$ of $\mathbf a_I$.
\begin{align}
 a_I^i&=\mathbf e_i\cdot\mathbf a
  =a_B^j\,\mathbf e_i\cdot\mathbf b_j=R_{ij}a_B^j,\notag\\
 a_B^j&=\mathbf b_j\cdot\mathbf a
  =a_I^i\,\mathbf b_j\cdot\mathbf e_i=R_{ij}a_I^i.
  \label{gnc:dcm-components}
\end{align}
The first line is Eq.~\eqref{gnc:configuration-map}: $R$ maps body
components to inertial components of the same vector,
\[
 R\colon\mathbb R^3\to\mathbb R^3,\qquad
 (a_B^1,a_B^2,a_B^3)\mapsto(a_I^1,a_I^2,a_I^3),\quad
 a_I^i=R_{ij}a_B^j .
\]
This is a \emph{passive rotation}: it changes the basis in which
$\mathbf a$ is described; $\mathbf a$ itself does not move.
For the \emph{active} interpretation, hold the inertial basis fixed and
let $\mathcal R$ rotate $\mathbf e_j$ into $\mathbf b_j$:
\begin{equation}\label{gnc:dcm-active}
 \mathcal R\mathbf e_j=\mathbf b_j,\qquad
 \mathcal R\mathbf a=a_I^j\mathbf b_j
 =R_{ij}a_I^j\mathbf e_i .
\end{equation}
Here the vector changes and both its input and output components use the
inertial basis. Goldstein's $A=R^T$ instead converts inertial components
to body components passively, or represents $\mathcal R^{-1}$ actively
in the fixed inertial basis (\S4.2, Eqs.~(4.18)--(4.19), pp.~142--143
\cite{dgdt:GPS2001}).

Orthonormality of each basis constrains the nine $R_{ij}$. Dot the
expansions~\eqref{gnc:dcm-basis} pairwise:
\begin{align}
 \delta_{jk}=\mathbf b_j\cdot\mathbf b_k
 &=R_{ij}R_{lk}\,\mathbf e_i\cdot\mathbf e_l
  =R_{ij}R_{lk}\delta_{il}=R_{ij}R_{ik},\notag\\
 \delta_{il}=\mathbf e_i\cdot\mathbf e_l
 &=R_{ij}R_{lk}\,\mathbf b_j\cdot\mathbf b_k
  =R_{ij}R_{lk}\delta_{jk}=R_{ij}R_{lj}.\label{gnc:dcm-orthogonality}
\end{align}
The first line is unchanged under $j\leftrightarrow k$, so it gives six
equations ($j\le k$) on nine entries; the regular-value proof of
Section~\ref{quat:sec:orthogonal-submanifold} makes the count $9-6=3$
rigorous. The nine direction cosines are therefore not independent
generalized coordinates.

Right-handedness, $\mathbf b_1\times\mathbf b_2=\mathbf b_3$, fixes the
determinant. In right-handed inertial components,
$(\mathbf u\times\mathbf v)\cdot\mathbf w=\epsilon_{ikl}u^iv^kw^l$, and
by Eq.~\eqref{gnc:dcm-basis} $\mathbf b_j$ has components $R_{ij}$:
\begin{equation}\label{gnc:dcm-determinant}
 \det R=\epsilon_{ikl}R_{i1}R_{k2}R_{l3}
 =(\mathbf b_1\times\mathbf b_2)\cdot\mathbf b_3
 =\mathbf b_3\cdot\mathbf b_3=1 .
\end{equation}
Hence $R\in GL(3,\mathbb R)$, the invertible real $3\times3$ matrices:
$\det R=1\ne0$, and the second line of Eqs.~\eqref{gnc:dcm-components}
inverts the first. Substituting $a_I^i=R_{ik}a_B^k$ and using the first
line of Eqs.~\eqref{gnc:dcm-orthogonality},
\[
 R_{ij}a_I^i=R_{ij}R_{ik}a_B^k=\delta_{jk}a_B^k=a_B^j .
\]
Equations~\eqref{gnc:dcm-orthogonality} and~\eqref{gnc:dcm-determinant}
place $R$ in $SO(3)\subset O(3)\subset GL(3,\mathbb R)$.

Engineering texts call such a matrix the direction cosine matrix (DCM).
The name does not fix which matrix is meant: $R^T$ also consists of
direction cosines, and many texts use it. Read the convention from the
defining equation. In the table, the space (inertial) basis is listed
first and the body basis second.
\begin{center}
\small
\begin{tabular}{@{}llll@{}}
 source & bases & direction cosines & matrix\\\hline
 ours & $\mathbf e_i$; $\mathbf b_j$ &
   $R_{ij}=\mathbf e_i\cdot\mathbf b_j$ & $R$\\
 Goldstein \cite{dgdt:GPS2001} & $\mathbf i,\mathbf j,\mathbf k$;
   $\mathbf i',\mathbf j',\mathbf k'$ &
   $a_{ij}=\cos\theta_{ij}$ & $A=R^T$\\
 Landau \cite{LLandauELifshitz1976} & $X,Y,Z$; $x_1,x_2,x_3$ &
   $n_{3x},n_{3y},n_{3z}$ & column 3 of $R$\\
 Wie \cite{BWie2008} & $\vec a_i$; $\vec b_i$ &
   $C_{ij}=\vec b_i\cdot\vec a_j$ & $C^{B/A}=R^T$\\
 Shuster \cite{MDShuster1993} & reference; body &
   attitude matrix & $A=R^T$\\
 Groves \cite{PDGroves2013} & $i$; $b$ &
   $\mathbf x^i=C^i_b\mathbf x^b$ & $C^i_b=R$
\end{tabular}
\end{center}

Goldstein (\S4.1, pp.~137--138; \S4.2, pp.~139--140) takes the primed
axes fixed in the body, so $\mathbf i,\mathbf j,\mathbf k$ are
$\mathbf e_1,\mathbf e_2,\mathbf e_3$ and $\mathbf i',\mathbf j',\mathbf k'$
are $\mathbf b_1,\mathbf b_2,\mathbf b_3$. In his Eq.~(4.2) the first index
of $\theta_{ij}$ is primed and the second unprimed, so
$\cos\theta_{ij}=\mathbf b_i\cdot\mathbf e_j=R_{ji}$, and his
$a_{ij}=\cos\theta_{ij}$ (Eq.~(4.11)) gives $A=R^T$. His
Eqs.~(4.5) and~(4.12$'$), $x'_i=a_{ij}x_j$, are the second of
Eqs.~\eqref{gnc:dcm-components}. His Eq.~(4.9),
$\sum_l\cos\theta_{lm'}\cos\theta_{lm}=\delta_{m'm}$, becomes
$R_{m'l}R_{ml}=\delta_{m'm}$ on substituting $\cos\theta_{lm}=R_{ml}$; his
Eq.~(4.15), $a_{ij}a_{ik}=\delta_{jk}$, becomes
$R_{ji}R_{ki}=\delta_{jk}$ on substituting $a_{ij}=R_{ji}$. Renaming
$(m',l,m)\to(i,j,l)$ in the first and $(j,i,k)\to(i,j,l)$ in the second
gives $R_{ij}R_{lj}=\delta_{il}$, the second of
Eqs.~\eqref{gnc:dcm-orthogonality}.

Landau (\S31, p.~96) uses fixed axes $X,Y,Z$ and body axes $x_1,x_2,x_3$,
and fixes the orientation by three independent angles, the Eulerian
angles of \S35 (p.~110). He does not name the direction cosines. In \S37
(p.~121) his $n_{3x},n_{3y},n_{3z}$, the fixed-axis components of the unit
vector $\mathbf n_3$ along $x_3$, are $R_{13},R_{23},R_{33}$.

Wie (\S5.1, pp.~323--325) writes Eq.~(5.1) as $\vec b_i=C_{ij}\vec a_j$
with $C_{ij}=\vec b_i\cdot\vec a_j$, calls $C^{B/A}$ the direction cosine
matrix ``to B from A'' (Eq.~(5.2)), and transforms components by
$\mathbf H'=C^{B/A}\mathbf H$ (Eq.~(5.8)). With $A$ the inertial frame,
$C^{B/A}=R^T$. His elementary rotations of Eq.~(5.9) are therefore
transposes of active ones: $C_k(\theta)=R_k(\theta)^T$, with
$R_1=R_x$, $R_2=R_y$, $R_3=R_z$. For example,
\[
 C_3(\theta)=
 \begin{pmatrix}\cos\theta&\sin\theta&0\\-\sin\theta&\cos\theta&0\\0&0&1
 \end{pmatrix}
 =R_z(\theta)^T .
\]
Transposing Eq.~\eqref{gnc:zyx-angles} gives the corresponding
yaw--pitch--roll product,
$R^T=R_x(\phi)^TR_y(\theta)^TR_z(\psi)^T=C_1(\phi)C_2(\theta)C_3(\psi)$.

Shuster's attitude matrix maps reference-frame components to body
components, $\mathbf a_B=A\mathbf a_I$, so $A=R^T$. Inertial-navigation
texts such as Groves write $C^\alpha_\beta$ for the matrix taking
$\beta$-frame components to $\alpha$-frame components, so the
body-to-inertial $C^i_b$ is $R$.

The components of a fixed vector turn opposite to the frame. Let the body
frame turn by $\alpha$ about $\mathbf e_3$:
\[
 \mathbf b_1=\cos\alpha\,\mathbf e_1+\sin\alpha\,\mathbf e_2,\qquad
 \mathbf b_2=-\sin\alpha\,\mathbf e_1+\cos\alpha\,\mathbf e_2,\qquad
 \mathbf b_3=\mathbf e_3 .
\]
For $\mathbf a=\mathbf e_1$, $a_I^i=\delta_{i1}$, and the second of
Eqs.~\eqref{gnc:dcm-components} with Eq.~\eqref{gnc:direction-cosines}
gives
\[
 a_B^j=R_{ij}\delta_{i1}=R_{1j}=\mathbf e_1\cdot\mathbf b_j:\qquad
 a_B^1=\cos\alpha,\quad a_B^2=-\sin\alpha,\quad a_B^3=0 .
\]
These are the components of $\mathbf e_1$ rotated by $-\alpha$. At
$\alpha=\pi/2$, $a_B=(0,-1,0)$, which checks against
$\mathbf b_2=-\mathbf e_1$.

\subsection{Coordinate changes, conjugation, and the adjoint representation}
\label{gnc:sec:conjugation}
Goldstein \S4.3, Eqs.~(4.40)--(4.41), pp.~148--149
\cite{dgdt:GPS2001}, transforms both the input and output coordinates of
an operator. Write the coordinate change as
\begin{equation}\label{gnc:coordinate-change}
 \begin{aligned}
 B&\colon\mathbb R^3_{\mathrm{old}}\longrightarrow
          \mathbb R^3_{\mathrm{new}},\qquad B\in GL(3,\mathbb R),\\
 x&\longmapsto y=Bx,\qquad y^i=B_{ij}x^j .
 \end{aligned}
\end{equation}
The subscripts label coordinate copies of $\mathbb R^3$ describing the
same vectors. For our body-to-inertial conversion, $B=R\in SO(3)$.
Let $A\in\operatorname{End}(\mathbb R^3_{\mathrm{old}})$ be any linear
map and $F\in\mathbb R^3_{\mathrm{old}}$ any input column vector; put
$G=AF$. Here $\operatorname{End}(V)$ is the vector space of linear maps
$V\to V$; in these coordinates $A$ is a $3\times3$ matrix.
Set $F'=BF$ and $G'=BG$. Then
\begin{equation}\label{gnc:operator-similarity}
 \begin{aligned}
 G'&=BG=BAF=BAB^{-1}F'=A'F',\\
 A'&=BAB^{-1},\qquad A\in\operatorname{End}(\mathbb R^3_{\mathrm{old}}),
 \quad A'\in\operatorname{End}(\mathbb R^3_{\mathrm{new}}).
 \end{aligned}
\end{equation}
This is Goldstein's \emph{similarity transformation}: the geometric
operator is unchanged. The factor $B^{-1}$ converts the input to old
coordinates; $A$ acts; $B$ converts the output to new coordinates.

For a matrix Lie group $G\subseteq GL(n,\mathbb R)$ with Lie algebra
$\mathfrak g=T_I G$ and $B\in G$, conjugation and its derivative at the identity are
\begin{equation}\label{gnc:conjugation-adjoint}
 \begin{aligned}
 c_B&\colon G\longrightarrow G,
       &h&\longmapsto BhB^{-1},\\
 \operatorname{Ad}_B=(dc_B)_I&\colon\mathfrak g\longrightarrow\mathfrak g,
       &X&\longmapsto
       \left.\frac{d}{dt}\bigl(B\exp(tX)B^{-1}\bigr)\right|_{t=0}
       =BXB^{-1}.
 \end{aligned}
\end{equation}
Thus each $\operatorname{Ad}_B$ is a similarity transformation restricted
to $\mathfrak g$. The \emph{adjoint representation} is the homomorphism
\[
 \operatorname{Ad}\colon G\longrightarrow GL(\mathfrak g),\qquad
 B\longmapsto\operatorname{Ad}_B,
\]
proved in Proposition~\ref{quat:prop:Ad}.
For $B\in SO(3)$, the hat map~\eqref{quat:hat} makes this concrete:
\begin{equation}\label{gnc:so3-adjoint-hat}
 \operatorname{Ad}_B\colon\mathfrak{so}(3)\longrightarrow\mathfrak{so}(3),
 \qquad\widehat{\omega}\longmapsto
 B\widehat{\omega}B^T=\widehat{B\omega}.
\end{equation}
Indeed, for every $v\in\mathbb R^3$, rotation invariance of the cross
product gives
\[
 (B\widehat{\omega}B^T)v
 =B(\omega\times B^Tv)
 =(B\omega)\times(BB^Tv)
 =(B\omega)\times v=\widehat{B\omega}\,v.
\]
With $B=R$, this identifies the coordinate laws
$\omega_I=R\omega_B$ and
$\widehat{\omega_I}=R\widehat{\omega_B}R^T$.

\subsection{Euler angles and Tait--Bryan coordinates}
Our convention remains $a_I=Ra_B$: a passive change of components.
The same matrix gives the active attitude of a body initially aligned
with the inertial axes. Put $c(\phi)\equiv\cos\phi$ and
$s(\phi)\equiv\sin\phi$. The elementary positive right-handed active
rotations are
\begin{equation}\label{gnc:elementary-rotations}
 \begin{aligned}
 R_x(\phi)&=\begin{pmatrix}1&0&0\\0&c(\phi)&-s(\phi)\\0&s(\phi)&c(\phi)\end{pmatrix},\\
 R_y(\phi)&=\begin{pmatrix}c(\phi)&0&s(\phi)\\0&1&0\\-s(\phi)&0&c(\phi)\end{pmatrix},\\
 R_z(\phi)&=\begin{pmatrix}c(\phi)&-s(\phi)&0\\s(\phi)&c(\phi)&0\\0&0&1\end{pmatrix}.
 \end{aligned}
\end{equation}

\paragraph{Goldstein's moving-axis $ZXZ$ sequence.}
The dictionary to Goldstein \S4.4, pp.~151--154
\cite{dgdt:GPS2001}, is
\[
 (\phi_G,\theta_G,\psi_G)=(\phi,\theta,\psi),\qquad
 (x,x')=(a_I,a_B),\qquad A_G=R_E^T=R_E^{-1}.
\]
Let $Q_k$ have the $k$th frame's basis vectors as columns in inertial
coordinates. His successive positive axis rotations give
\begin{equation}\label{gnc:zxz-frames}
 \begin{aligned}
 Q_0&=I,& Q_1&=R_z(\phi),\\
 Q_2&=Q_1R_x(\theta),& Q_3&=Q_2R_z(\psi)=R_E.
 \end{aligned}
\end{equation}
The physical axes are $e_3$, $Q_1e_1$ (the line of nodes), and $Q_2e_3$.
\begin{center}
\begingroup
\begin{tikzpicture}[
 x={(0.34cm,-0.54cm)},y={(0.94cm,0.20cm)},z={(0cm,0.82cm)},
 >=Stealth,font=\small,
 before/.style={->,gray!65,dashed},
 after/.style={->,blue!70!black,thick},
 axis/.style={->,orange!85!black,very thick},
 turn/.style={->,red!70!black,thick}]
\pgfmathsetmacro{\ca}{cos(35)}
\pgfmathsetmacro{\sa}{sin(35)}
\pgfmathsetmacro{\cb}{cos(40)}
\pgfmathsetmacro{\sb}{sin(40)}
\pgfmathsetmacro{\cg}{cos(30)}
\pgfmathsetmacro{\sg}{sin(30)}
\begin{scope}
 \draw[before] (0,0,0)--(1.65,0,0) node[right] {$e_1$};
 \draw[before] (0,0,0)--(0,1.65,0) node[right] {$e_2$};
 \draw[after] (0,0,0)--({1.65*\ca},{1.65*\sa},0) node[right] {$\xi$};
 \draw[after] (0,0,0)--({-1.65*\sa},{1.65*\ca},0) node[above] {$\eta$};
 \draw[axis] (0,0,0)--(0,0,1.65) node[above] {$e_3=\zeta$};
 \draw[turn] plot[domain=0:35,samples=16,variable=\t]
  ({0.85*cos(\t)},{0.85*sin(\t)},0);
 \node at ({1.15*cos(17.5)},{1.15*sin(17.5)},0) {$\phi$};
 \node at (0,0,-1.8) {$Q_0\longrightarrow Q_1$};
\end{scope}
\begin{scope}[xshift=4.2cm]
 \draw[before] (0,0,0)--({-1.65*\sa},{1.65*\ca},0);
 \draw[before] (0,0,0)--(0,0,1.65);
 \draw[after] (0,0,0)--({-1.65*\sa*\cb},{1.65*\ca*\cb},{1.65*\sb})
  node[above] {$\eta'$};
 \draw[after] (0,0,0)--({1.65*\sa*\sb},{-1.65*\ca*\sb},{1.65*\cb})
  node[left] {$\zeta'$};
 \draw[axis] (0,0,0)--({1.65*\ca},{1.65*\sa},0)
  node[right] {$\xi=\xi'$};
 \draw[turn] plot[domain=0:40,samples=16,variable=\t]
  ({-0.85*\sa*cos(\t)},{0.85*\ca*cos(\t)},{0.85*sin(\t)});
 \node at ({-1.1*\sa*cos(20)},{1.1*\ca*cos(20)},{1.1*sin(20)}) {$\theta$};
 \node at (0,0,-1.8) {$Q_1\longrightarrow Q_2$};
\end{scope}
\begin{scope}[xshift=8.4cm]
 \draw[before] (0,0,0)--({1.65*\ca},{1.65*\sa},0);
 \draw[before] (0,0,0)--({-1.65*\sa*\cb},{1.65*\ca*\cb},{1.65*\sb});
 \draw[after] (0,0,0)--({1.65*(\ca*\cg-\sa*\cb*\sg)},
  {1.65*(\sa*\cg+\ca*\cb*\sg)},{1.65*\sb*\sg}) node[right] {$b_1$};
 \draw[after] (0,0,0)--({1.65*(-\ca*\sg-\sa*\cb*\cg)},
  {1.65*(-\sa*\sg+\ca*\cb*\cg)},{1.65*\sb*\cg}) node[above] {$b_2$};
 \draw[axis] (0,0,0)--({1.65*\sa*\sb},{-1.65*\ca*\sb},{1.65*\cb})
  node[left] {$\zeta'=b_3$};
 \draw[turn] plot[domain=0:30,samples=16,variable=\t]
  ({0.85*(\ca*cos(\t)-\sa*\cb*sin(\t))},
   {0.85*(\sa*cos(\t)+\ca*\cb*sin(\t))},{0.85*\sb*sin(\t)});
 \node[below] at ({1.2*(\ca*cos(15)-\sa*\cb*sin(15))},
  {1.2*(\sa*cos(15)+\ca*\cb*sin(15))},{1.2*\sb*sin(15)}) {$\psi$};
 \node at (0,0,-1.8) {$Q_2\longrightarrow Q_3$};
\end{scope}
\end{tikzpicture}
\par\small $ZXZ$: previous axes dashed, new axes blue, rotation axis orange.
Illustrated at $(\phi,\theta,\psi)=(35^\circ,40^\circ,30^\circ)$.
\endgroup
\end{center}

For one fixed vector $a=a_k^jQ_ke_j$, $a_0=a_I$ and $a_3=a_B$:
\begin{equation}\label{gnc:goldstein-passive-sequence}
 \begin{aligned}
 a_1&=\underbrace{R_z(-\phi)}_{D_G}a_0,\\
 a_2&=\underbrace{R_x(-\theta)}_{C_G}a_1,\\
 a_3&=\underbrace{R_z(-\psi)}_{B_G}a_2
     =B_GC_GD_Ga_I=A_Ga_I,\\
 A_G&=R_z(-\psi)R_x(-\theta)R_z(-\phi).
 \end{aligned}
\end{equation}
Thus, in our body-to-inertial convention,
\begin{equation}\label{gnc:zxz-angles}
 R_E=R_z(\phi)R_x(\theta)R_z(\psi).
\end{equation}
The multiplication is explicit:
\begin{align}
 R_z(\phi)R_x(\theta)
 &=\begin{pmatrix}
 c(\phi)&-s(\phi)c(\theta)&s(\phi)s(\theta)\\
 s(\phi)&c(\phi)c(\theta)&-c(\phi)s(\theta)\\
 0&s(\theta)&c(\theta)
 \end{pmatrix},\notag\\
 R_E&=\begin{pmatrix}
 c(\phi)c(\psi)-s(\phi)c(\theta)s(\psi)&
 -c(\phi)s(\psi)-s(\phi)c(\theta)c(\psi)&s(\phi)s(\theta)\\
 s(\phi)c(\psi)+c(\phi)c(\theta)s(\psi)&
 -s(\phi)s(\psi)+c(\phi)c(\theta)c(\psi)&-c(\phi)s(\theta)\\
 s(\theta)s(\psi)&s(\theta)c(\psi)&c(\theta)
 \end{pmatrix}.\label{gnc:zxz-expanded}
\end{align}
Its transpose is Goldstein's Eq.~(4.46).

\paragraph{Active readings: specify the axes.}
The positive physical rotations about those moving axes, all represented
in the fixed inertial basis, are
\[
 U_1=Q_1,\qquad U_2=Q_1R_x(\theta)Q_1^T,\qquad
 U_3=Q_2R_z(\psi)Q_2^T.
\]
Their chronological product is
\[
 U_2U_1=Q_1R_x(\theta)Q_1^TQ_1=Q_2,\qquad
 U_3U_2U_1=Q_2R_z(\psi)Q_2^TQ_2=R_E.
\]
Equivalently, the same $R_E$ actively rotates vectors about fixed axes in
chronological order $z(\psi),x(\theta),z(\phi)$.
Reading Goldstein's $D_G,C_G,B_G$ actively instead gives
\[
 v_I\xmapsto{D_G}v_I^{(1)}\xmapsto{C_G}v_I^{(2)}
 \xmapsto{B_G}v_I^{(3)}=A_Gv_I=R_E^{-1}v_I,
\]
with all components in the inertial basis and fixed-axis rotations
$z(-\phi),x(-\theta),z(-\psi)$.
Goldstein defines the angles by rotating axes and uses the passive
component law. For physical body motion the active $R_E$ is equally
useful; its inverse describes the opposite motion.

\paragraph{Goldstein's $y$-convention: $ZYZ$.}
Replace the second moving axis by $Q_1e_2$:
\begin{equation}\label{gnc:zyz-angles}
 R_{E,y}(\phi,\theta,\psi)=R_z(\phi)R_y(\theta)R_z(\psi),\qquad
 A_{G,y}=R_{E,y}^T.
\end{equation}
The products are
\begin{align}
 R_z(\phi)R_y(\theta)
 &=\begin{pmatrix}
 c(\phi)c(\theta)&-s(\phi)&c(\phi)s(\theta)\\
 s(\phi)c(\theta)&c(\phi)&s(\phi)s(\theta)\\
 -s(\theta)&0&c(\theta)
 \end{pmatrix},\notag\\
 R_{E,y}&=\begin{pmatrix}
 c(\phi)c(\theta)c(\psi)-s(\phi)s(\psi)&
 -c(\phi)c(\theta)s(\psi)-s(\phi)c(\psi)&c(\phi)s(\theta)\\
 s(\phi)c(\theta)c(\psi)+c(\phi)s(\psi)&
 -s(\phi)c(\theta)s(\psi)+c(\phi)c(\psi)&s(\phi)s(\theta)\\
 -s(\theta)c(\psi)&s(\theta)s(\psi)&c(\theta)
 \end{pmatrix}.\label{gnc:zyz-expanded}
\end{align}
The exact dictionary between his $x$- and $y$-conventions follows from
\begin{align*}
 R_x(\theta)&=R_z(-\pi/2)R_y(\theta)R_z(\pi/2),\\
 R_E(\phi,\theta,\psi)
 &=R_z(\phi-\pi/2)R_y(\theta)R_z(\psi+\pi/2)\\
 &=R_{E,y}(\phi-\pi/2,\theta,\psi+\pi/2).
\end{align*}
The active/passive rules above apply with $x$ replaced by $y$.

\paragraph{Standard $ZYX$ yaw--pitch--roll convention.}
We use the standard aerospace $3$--$2$--$1$ sequence with yaw $\psi$,
pitch $\theta$, and roll $\phi$; our matrix direction is body to inertial.
For the same angles, MathWorks's \texttt{eul2rotm} returns $R$, while its
aerospace \texttt{angle2dcm} returns $R^T$
\cite{MathWorksRotationConventions}.
The positive moving-axis
sequence is $z(\psi),y'(\theta),x''(\phi)$; equivalently, the fixed-axis
active order is $x(\phi),y(\theta),z(\psi)$:
\begin{equation}\label{gnc:zyx-angles}
 R=R_z(\psi)R_y(\theta)R_x(\phi),\qquad
 a_B=R_x(-\phi)R_y(-\theta)R_z(-\psi)a_I.
\end{equation}
The corresponding multiplication gives
\begin{align}
 R_y(\theta)R_x(\phi)
 &=\begin{pmatrix}
 c(\theta)&s(\theta)s(\phi)&s(\theta)c(\phi)\\
 0&c(\phi)&-s(\phi)\\
 -s(\theta)&c(\theta)s(\phi)&c(\theta)c(\phi)
 \end{pmatrix},\notag\\
 R&=\begin{pmatrix}
 c(\psi)c(\theta)&c(\psi)s(\theta)s(\phi)-s(\psi)c(\phi)&
 c(\psi)s(\theta)c(\phi)+s(\psi)s(\phi)\\
 s(\psi)c(\theta)&s(\psi)s(\theta)s(\phi)+c(\psi)c(\phi)&
 s(\psi)s(\theta)c(\phi)-c(\psi)s(\phi)\\
 -s(\theta)&c(\theta)s(\phi)&c(\theta)c(\phi)
 \end{pmatrix}.\label{gnc:zyx-expanded}
\end{align}
Goldstein's ``$xyz$-convention'' means three distinct axes; here we fix
its order as $ZYX$. The dictionary of chronological rotations is
\begin{center}
\begin{tabular}{@{}lll@{}}
 convention & moving axes (intrinsic) & fixed axes (active)\\\hline
 $x$: $ZXZ$ & $z(\phi),x'(\theta),z''(\psi)$ & $z(\psi),x(\theta),z(\phi)$\\
 $y$: $ZYZ$ & $z(\phi),y'(\theta),z''(\psi)$ & $z(\psi),y(\theta),z(\phi)$\\
 ours: $ZYX$ & $z(\psi),y'(\theta),x''(\phi)$ & $x(\phi),y(\theta),z(\psi)$
\end{tabular}
\end{center}
Goldstein's names are yaw (heading), pitch (attitude), and roll (bank). For the same attitude,
\[
 R_E(\phi_x,\theta_x,\psi_x)
 =R_{E,y}(\phi_y,\theta_y,\psi_y)
 =R_z(\psi)R_y(\theta)R_x(\phi)=R.
\]
Use the entries of this common matrix to convert between all three:
\begin{equation}\label{gnc:euler-angle-dictionary}
 \begin{aligned}
 \phi_x&=\operatorname{atan2}(R_{13},-R_{23}),&
 \theta_x&=\arccos R_{33},&
 \psi_x&=\operatorname{atan2}(R_{31},R_{32}),\\
 \phi_y&=\operatorname{atan2}(R_{23},R_{13}),&
 \theta_y&=\arccos R_{33},&
 \psi_y&=\operatorname{atan2}(R_{32},-R_{31}),\\
 \psi&=\operatorname{atan2}(R_{21},R_{11}),&
 \theta&=\arcsin(-R_{31}),&
 \phi&=\operatorname{atan2}(R_{32},R_{33}).
 \end{aligned}
\end{equation}
Here $0<\theta_x,\theta_y<\pi$ and $-\pi/2<\theta<\pi/2$; outer angles
are taken modulo $2\pi$. In particular,
\[
 \theta=\arcsin\bigl(-s(\theta_x)s(\psi_x)\bigr)
       =\arcsin\bigl(s(\theta_y)c(\psi_y)\bigr).
\]
Both repeated-axis conventions are singular at $\sin\theta=0$; for
example $R_E(\phi,0,\psi)=R_z(\phi+\psi)$.
Our Tait--Bryan angles are regular near the identity and singular at
$\cos\theta=0$.

For the rate calculation use our $ZYX$ choice~\eqref{gnc:zyx-angles}.
Let $E_i=\widehat{\mathbf e_i}$ be the generators of
Eq.~\eqref{quat:frankel-Ei}. Differentiating and multiplying by $R^T$ gives
\begin{align*}
 R^T\dot R
 &=E_1\dot\phi+R_x^TE_2R_x\dot\theta\\
 &\quad+R_x^TR_y^TE_3R_yR_x\dot\psi.
\end{align*}
Define $\widehat{\boldsymbol\omega}_B=R^T\dot R$.
Using $R^T\widehat{\mathbf a}R=\widehat{R^T\mathbf a}$ gives
\begin{equation}\label{gnc:zyx-rate-map}
 \boldsymbol\omega_B=
 \begin{pmatrix}
  1&0&-s(\theta)\\
  0&c(\phi)&s(\phi)c(\theta)\\
  0&-s(\phi)&c(\phi)c(\theta)
 \end{pmatrix}
 \begin{pmatrix}\dot\phi\\\dot\theta\\\dot\psi\end{pmatrix}.
\end{equation}
The determinant of the rate matrix is
$c(\theta)\bigl(c(\phi)^2+s(\phi)^2\bigr)=c(\theta)$.
When $c(\theta)\ne0$, the inverse function theorem gives a local
coordinate chart after restricting the angle ranges. At
$\theta=\pm\pi/2$ the angle map loses rank (gimbal lock); $SO(3)$ remains
smooth. A unit quaternion avoids this angle singularity, but uses four
constrained components and the two representatives $q,-q$ of one rotation.

\subsection{Cayley--Klein parameters, quaternions, and angles}
\label{gnc:sec:cayley-klein}
Use subscripts to distinguish Goldstein's complex parameters
(\S4.5, pp.~154--155 \cite{dgdt:GPS2001}) from the angles:
\begin{equation}\label{gnc:ck-constraints}
 \begin{gathered}
 \alpha_{\mathrm{CK}},\beta_{\mathrm{CK}},\gamma_{\mathrm{CK}},
 \delta_{\mathrm{CK}}\in\mathbb C,\qquad
 \gamma_{\mathrm{CK}}=-\beta_{\mathrm{CK}}^*,\quad
 \delta_{\mathrm{CK}}=\alpha_{\mathrm{CK}}^*,\\
 |\alpha_{\mathrm{CK}}|^2+|\beta_{\mathrm{CK}}|^2=1,\qquad
 K=\begin{pmatrix}\alpha_{\mathrm{CK}}&\beta_{\mathrm{CK}}\\
                  -\beta_{\mathrm{CK}}^*&\alpha_{\mathrm{CK}}^*\end{pmatrix}.
 \end{gathered}
\end{equation}
Here $^*$ denotes complex conjugation. The constraints give
\[
 K^\dagger K=(|\alpha_{\mathrm{CK}}|^2+|\beta_{\mathrm{CK}}|^2)I=I,
 \qquad \det K=|\alpha_{\mathrm{CK}}|^2+|\beta_{\mathrm{CK}}|^2=1;
 \qquad K\in SU(2).
\]

\paragraph{Dictionary to our quaternion.}
Goldstein's four real \emph{Euler parameters} are
\begin{equation}\label{gnc:euler-parameters}
 \begin{gathered}
 (e_0,e_1,e_2,e_3)_G=(q_0,q^1,q^2,q^3)\in\mathbb R^4,\qquad
 q_0^2+(q^1)^2+(q^2)^2+(q^3)^2=1,\\
 \alpha_{\mathrm{CK}}=q_0+iq^3,\qquad
 \beta_{\mathrm{CK}}=q^2+iq^1,\qquad
 q=q_0+q^je_j\in S^3\subset\mathbb H.
 \end{gathered}
\end{equation}
The $e_j$ in $q$ are our Hamilton units~\eqref{quat:basis}; the $i$ in
$\alpha_{\mathrm{CK}},\beta_{\mathrm{CK}}$ is the complex unit.
Conversely,
\[
 q_0=\Re\alpha_{\mathrm{CK}},\quad q^3=\Im\alpha_{\mathrm{CK}},\qquad
 q^2=\Re\beta_{\mathrm{CK}},\quad q^1=\Im\beta_{\mathrm{CK}}.
\]
Under our quaternion--Pauli dictionary~\eqref{quat:dictionary},
\begin{equation}\label{gnc:ck-su2-dictionary}
 U(q)=q_0I-iq^j\sigma_j
 =\begin{pmatrix}q_0-iq^3&-q^2-iq^1\\q^2-iq^1&q_0+iq^3\end{pmatrix},
 \qquad K=U(q)^\dagger=U(\bar q).
\end{equation}
Thus the covering map~\eqref{quat:Phi} gives
\begin{equation}\label{gnc:ck-rotation}
 \begin{aligned}
 R(q)&=\Phi(U(q))
 =(q_0^2-\mathbf q\cdot\mathbf q)I
   +2\mathbf q\mathbf q^T+2q_0\widehat{\mathbf q},\\
 A_G&=\Phi(K)=\Phi(U(q)^{-1})=R(q)^{-1}=R(q)^T,
 \qquad\mathbf q=(q^1,q^2,q^3)^T.
 \end{aligned}
\end{equation}
This is Eq.~\eqref{quat:euler-rodrigues}. For example, Goldstein's complex
entry becomes
\begin{align*}
 (A_G)_{13}
 &=\gamma_{\mathrm{CK}}\delta_{\mathrm{CK}}
   -\alpha_{\mathrm{CK}}\beta_{\mathrm{CK}}\\
 &=(-q^2+iq^1)(q_0-iq^3)-(q_0+iq^3)(q^2+iq^1)\\
 &=2(q^1q^3-q_0q^2)=R(q)_{31}.
\end{align*}
All entries are real quadratic expressions; $q$ and $-q$ give the same
$R\in SO(3)$. These are four constrained real coordinates for attitude,
with no Euler-angle singularity.

\paragraph{Angles to quaternions and back.}
Put $q_a(\phi)=c(\phi/2)+e_a s(\phi/2)$ for $a=1,2,3$. Hamilton multiplication
follows the same order as the body-to-inertial matrices:
\begin{equation}\label{gnc:angle-quaternion-products}
 \begin{aligned}
 q_{ZXZ}&=q_3(\phi)q_1(\theta)q_3(\psi),\\
 q_{ZYZ}&=q_3(\phi)q_2(\theta)q_3(\psi),\\
 q_{ZYX}&=q_3(\psi)q_2(\theta)q_1(\phi).
 \end{aligned}
\end{equation}
Writing the results scalar first, with $u=(\phi+\psi)/2$ and
$v=(\phi-\psi)/2$,
\begin{equation}\label{gnc:proper-euler-quaternions}
 q_{ZXZ}\ \longleftrightarrow\
 \begin{pmatrix}c(\theta/2)c(u)\\s(\theta/2)c(v)\\
 s(\theta/2)s(v)\\c(\theta/2)s(u)\end{pmatrix},\qquad
 q_{ZYZ}\ \longleftrightarrow\
 \begin{pmatrix}c(\theta/2)c(u)\\-s(\theta/2)s(v)\\
 s(\theta/2)c(v)\\c(\theta/2)s(u)\end{pmatrix}.
\end{equation}
For Tait--Bryan angles,
\begin{equation}\label{gnc:tait-bryan-quaternion}
 q_{ZYX}\ \longleftrightarrow\
 \begin{pmatrix}
 c(\psi/2)c(\theta/2)c(\phi/2)+s(\psi/2)s(\theta/2)s(\phi/2)\\
 c(\psi/2)c(\theta/2)s(\phi/2)-s(\psi/2)s(\theta/2)c(\phi/2)\\
 c(\psi/2)s(\theta/2)c(\phi/2)+s(\psi/2)c(\theta/2)s(\phi/2)\\
 s(\psi/2)c(\theta/2)c(\phi/2)-c(\psi/2)s(\theta/2)s(\phi/2)
 \end{pmatrix}.
\end{equation}
In particular Goldstein's $ZXZ$ angles give directly
\[
 \alpha_{\mathrm{CK}}=c(\theta/2)e^{iu},\qquad
 \beta_{\mathrm{CK}}=s(\theta/2)\bigl(s(v)+ic(v)\bigr)
                   =i\,s(\theta/2)e^{-iv}.
\]
For the reverse conversion, compute $R(q)$ from
Eq.~\eqref{gnc:ck-rotation} and apply
Eq.~\eqref{gnc:euler-angle-dictionary} on its stated angle branches.

\subsection{Euler's theorem: one axis and one finite angle}
\label{gnc:sec:euler-theorem}
Goldstein \S4.6, pp.~155--161, Eqs.~(4.48), (4.56)--(4.61)
\cite{dgdt:GPS2001}. In our notation the statement is
\[
 R\in SO(3)\quad\Longrightarrow\quad
 \exists\,n\in S^2:\ Rn=n.
\]
Indeed, putting $d=\det(R-I)$,
\begin{align*}
 (R-I)R^T&=I-R^T,\\
 d&=\det\bigl((R-I)R^T\bigr)
   =\det(I-R^T)=\det(I-R)=-\det(R-I)=-d,\\
 d&=0,\qquad \ker(R-I)\ne\{0\}.
\end{align*}
Choose $n\in\ker(R-I)$ with $|n|=1$, $|u|=1$, $u\perp n$, and
$v=n\times u$. Then $S=(u\ v\ n)\in SO(3)$ and
\begin{equation}\label{gnc:rotation-normal-form}
 S^TRS=\begin{pmatrix}c(\phi)&-s(\phi)&0\\s(\phi)&c(\phi)&0\\0&0&1\end{pmatrix}.
\end{equation}
The sign of $n$ can be chosen so that $0\le\phi\le\pi$.
Consequently,
\begin{equation}\label{gnc:rotation-spectrum}
 \begin{gathered}
 \operatorname{spec}_{\mathbb C}R=\{1,e^{i\phi},e^{-i\phi}\},\qquad
 \operatorname{tr}R=1+2c(\phi),\\
 \dim\ker(R-I)=\begin{cases}3,&R=I,\\1,&R\ne I.\end{cases}
 \end{gathered}
\end{equation}
The characteristic polynomial and Cayley--Hamilton identity are
\begin{equation}\label{gnc:rotation-cayley-hamilton}
 \begin{aligned}
 \chi_R(\lambda)&=(\lambda-1)(\lambda^2-2c(\phi)\lambda+1),\\
 \chi_R(R)&=R^3-\tau R^2+\tau R-I=0,\qquad\tau=\operatorname{tr}R.
 \end{aligned}
\end{equation}
The determinant argument uses odd dimension; $-I_2\in SO(2)$ has no
nonzero fixed vector.

\paragraph{Finite rotations and the exponential map.}
Goldstein \S4.7, pp.~161--163, Eqs.~(4.62)--(4.63), uses axis--angle
coordinates for the same rotation. Put $N=\widehat n\in\mathfrak{so}(3)$,
$Nx=n\times x$. The vector triple product gives
\[
 N^2x=n\times(n\times x)=n(n\cdot x)-x,\qquad
 N^2=nn^T-I,\qquad N^3=Nnn^T-N=-N.
\]
Reducing the exponential series with $N^3=-N$,
\begin{align}
 \exp(\phi N)
 &=I+\sum_{m=0}^{\infty}\frac{(-1)^m\phi^{2m+1}}{(2m+1)!}\,N
    +\sum_{m=1}^{\infty}\frac{(-1)^{m-1}\phi^{2m}}{(2m)!}\,N^2\notag\\
 &=I+s(\phi)N+(1-c(\phi))N^2\notag\\
 &=c(\phi)I+(1-c(\phi))nn^T+s(\phi)\widehat n
 =:R_{n,\phi}.\label{gnc:rodrigues-matrix}
\end{align}
Thus the exponential is surjective,
\[
 \exp\colon\mathfrak{so}(3)\longrightarrow SO(3),\qquad
 \phi\widehat n\longmapsto R_{n,\phi},
\]
and the corresponding active vector map is precisely
\begin{equation}\label{gnc:finite-vector-map}
 \begin{aligned}
 R_{n,\phi}&\colon\mathbb R^3\longrightarrow\mathbb R^3,\\
 x&\longmapsto c(\phi)x+(1-c(\phi))n(n\cdot x)+s(\phi)n\times x.
 \end{aligned}
\end{equation}
This agrees with Eq.~\eqref{quat:rotate-about-n}.
Goldstein's Eq.~(4.62) has the opposite active sense:
\begin{equation}\label{gnc:goldstein-finite-dictionary}
 \begin{aligned}
 A_G&=R_{n,\phi}^T=R_{n,-\phi},\\
 A_Gx&=c(\phi)x+(1-c(\phi))n(n\cdot x)+s(\phi)x\times n.
 \end{aligned}
\end{equation}

\paragraph{Recovering the axis and angle.}
Equations~\eqref{gnc:rotation-spectrum} and~\eqref{gnc:rodrigues-matrix}
give
\begin{equation}\label{gnc:axis-angle-recovery}
 \begin{gathered}
 \phi=\arccos\!\left(\frac{\operatorname{tr}R-1}{2}\right),\qquad
 R-R^T=2s(\phi)\widehat n,\\
 n=\frac1{2s(\phi)}\begin{pmatrix}R_{32}-R_{23}\\R_{13}-R_{31}\\R_{21}-R_{12}\end{pmatrix},
 \qquad 0<\phi<\pi.
 \end{gathered}
\end{equation}
At $R=I$, every axis is possible. At $\phi=\pi$,
\[
 \frac{R+I}{2}=nn^T,\qquad n\sim-n.
\]
For example, choose $j$ with maximal $R_{jj}$ and take
$n_j=\sqrt{(R_{jj}+1)/2}>0$; then
$n_i=(R_{ij}+\delta_{ij})/(2n_j)$.

\paragraph{One rotation, several representations.}
With convention subscripts on the Euler angles,
\begin{equation}\label{gnc:rotation-representation-dictionary}
 \begin{aligned}
 R&=R_z(\phi_E)R_x(\theta_E)R_z(\psi_E)\\
  &=R_z(\psi_T)R_y(\theta_T)R_x(\phi_T)\\
  &=\exp(\phi\widehat n)=\Phi(U(q)),\\
 q&=\pm\bigl(c(\phi/2)+n^je_j s(\phi/2)\bigr),\qquad
 K=U(q)^\dagger,\quad A_G=\Phi(K)=R^T.
 \end{aligned}
\end{equation}
Use Eqs.~\eqref{gnc:euler-angle-dictionary},
\eqref{gnc:proper-euler-quaternions} and~\eqref{gnc:tait-bryan-quaternion}
for the conversions. In particular, for the principal angle $0\le\phi\le\pi$,
\begin{equation}\label{gnc:goldstein-half-angle}
 c(\phi/2)=|q_0|
 =\left|c(\theta_E/2)c\bigl((\phi_E+\psi_E)/2\bigr)\right|.
\end{equation}
Goldstein's Eq.~(4.63) chooses the continuous lift near the identity;
the absolute value selects the principal angle globally.

For a body-fixed vector $a_B$, its finite displacement between two times is
\[
 a_I(t_2)=R(t_2)a_B=R(t_2)R(t_1)^Ta_I(t_1).
\]
Euler's theorem applies to $R(t_2)R(t_1)^T$; its axis generally differs
from the instantaneous angular-velocity direction.

\paragraph{Chasles: translation, rotation, and a screw axis.}
The free rigid displacement is
\[
 (r,R)\in SE(3),\qquad x\longmapsto r+Rx.
\]
For $R=R_{n,\phi}\ne I$, put $h=n\cdot r$, $r_\perp=r-hn$, and
\[
 p=\frac12\left(r_\perp+\cot(\phi/2)n\times r_\perp\right).
\]
Since $R=c(\phi)I+s(\phi)N$ on $n^\perp$ and $N^2=-I$ there,
\begin{align*}
 (I-R)p
 &=\tfrac12\bigl((1-c(\phi))I-s(\phi)N\bigr)
       \bigl(I+\cot(\phi/2)N\bigr)r_\perp\\
 &=\tfrac12\Bigl((1-c(\phi)+s(\phi)\cot(\phi/2))I
       +((1-c(\phi))\cot(\phi/2)-s(\phi))N\Bigr)r_\perp\\
 &=r_\perp.
\end{align*}
Hence
\[
 r+Rx=p+R(x-p)+hn.
\]
The axis is $p+\mathbb R n$ and the translation along it is $hn$.
For $R=I$ the motion is pure translation.


\section{Angular velocity and the rotating-frame derivative}
Let $R(t)$ map body components to inertial components, with
$R^{T}R=I$ and $\det R=1$. Differentiating the first identity gives
\begin{align}
 \dot R^{T}R+R^{T}\dot R&=0,\notag\\
 (R^{T}\dot R)^{T}&=-R^{T}\dot R.\label{gnc:skew}
\end{align}
Equation~\eqref{gnc:skew} makes $R^{T}\dot R$ a skew-symmetric matrix.
Define $\widehat{\boldsymbol\omega}_B=R^{T}\dot R$ by
$\widehat{\boldsymbol\omega}_B\mathbf a=\boldsymbol\omega_B\times\mathbf a$.
Multiplication by $R$ then yields
\begin{equation}\label{gnc:attitude-kinematics}
 \dot R=R\widehat{\boldsymbol\omega}_B.
\end{equation}
For any physical vector with components $\mathbf a_I=R\mathbf a_B$,
apply the product rule and substitute Eq.~\eqref{gnc:attitude-kinematics}:
\begin{align}
 \dot{\mathbf a}_I
 &=\dot R\mathbf a_B+R\dot{\mathbf a}_B\notag\\
 &=R\bigl(\boldsymbol\omega_B\times\mathbf a_B
                  +\dot{\mathbf a}_B\bigr).
 \label{gnc:transport}
\end{align}
The derivative on the left is taken in an inertial frame. The derivative of
body components on the right includes the change of basis through the
cross-product term.

\section{Momentum balances about the center of mass}
Assume a closed rigid body with constant mass $m>0$, constant positive-definite
body inertia $J_B$, and a body origin at its center of mass. Let
$\mathbf r_I$ and $\mathbf v_I$ be its inertial position and velocity.
The total external force is $\mathbf F_I$; the external moment about the
center of mass has components $\boldsymbol\tau_B$.
With linear momentum $\mathbf p_I=m\mathbf v_I$, Newton's balance is
\begin{equation}\label{gnc:linear-balance}
 \dot{\mathbf r}_I=\mathbf v_I,\qquad
 \dot{\mathbf p}_I=\mathbf F_I,\qquad
 m\dot{\mathbf v}_I=\mathbf F_I.
\end{equation}
Define angular momentum $\mathbf h_B=J_B\boldsymbol\omega_B$, so that
$\mathbf h_I=R\mathbf h_B$. Its inertial balance is
$\dot{\mathbf h}_I=R\boldsymbol\tau_B$. Substitution into
Eq.~\eqref{gnc:transport} gives
\begin{align}
 R\bigl(\dot{\mathbf h}_B+
        \boldsymbol\omega_B\times\mathbf h_B\bigr)
   &=R\boldsymbol\tau_B,\notag\\
 J_B\dot{\boldsymbol\omega}_B+
 \boldsymbol\omega_B\times(J_B\boldsymbol\omega_B)
   &=\boldsymbol\tau_B.\label{gnc:euler}
\end{align}
The second line uses constant $J_B$ and multiplication by $R^{T}$.
Similarly, setting $\mathbf v_I=R\mathbf v_B$ in
Eq.~\eqref{gnc:linear-balance}, and using Eq.~\eqref{gnc:transport}, gives
\begin{equation}\label{gnc:body-translation}
 m\bigl(\dot{\mathbf v}_B+
       \boldsymbol\omega_B\times\mathbf v_B\bigr)=\mathbf F_B,
 \qquad \mathbf F_B=R^{T}\mathbf F_I.
\end{equation}
Equations~\eqref{gnc:linear-balance}, \eqref{gnc:attitude-kinematics}, and
\eqref{gnc:euler} form a constant-mass rigid-body six-DOF model after the
force and moment laws are specified. Rocket mass loss requires open-system
momentum balances; replacing $m$ by $m(t)$ alone does not derive them.

For zero external moment, Eq.~\eqref{gnc:euler} implies
$\dot{\mathbf h}_I=0$, although $\mathbf h_B$ generally changes.
The rotational kinetic energy satisfies
\begin{align}
 T_{\mathrm{rot}}&=\tfrac12\boldsymbol\omega_B^{T}
                            J_B\boldsymbol\omega_B,\notag\\
 \dot T_{\mathrm{rot}}
 &=\boldsymbol\omega_B\cdot J_B\dot{\boldsymbol\omega}_B\notag\\
 &=\boldsymbol\omega_B\cdot\boldsymbol\tau_B
 -\boldsymbol\omega_B\cdot
       \bigl(\boldsymbol\omega_B\times J_B\boldsymbol\omega_B\bigr)\notag\\
 &=\boldsymbol\omega_B\cdot\boldsymbol\tau_B.
\end{align}
The cross-product term vanishes because it is perpendicular to
$\boldsymbol\omega_B$. Zero moment therefore conserves both inertial
angular momentum and rotational kinetic energy.

\section{A free six-DOF model and its numerical realization}
\label{gnc:sec:free-simulation}
Take the body origin at the center of mass, SI units, constant mass $m>0$
and constant physical inertia $J=J_B=J^T>0$. Its principal moments satisfy
$J_i\le J_j+J_k$. Write $\mathcal R:S^3\to SO(3)$ for our Hamilton,
scalar-first quaternion map, $q=(q_0,\mathbf q)$, with
$\mathcal R(q)\mathbf a_B=\mathbf a_I$. The state and evolution are
\begin{align}
 \mathcal X&=\mathbb R^3\times\mathbb R^3\times S^3\times\mathbb R^3,
 &z&=(\mathbf r_I,\mathbf v_I,q,\boldsymbol\omega_B),\notag\\
 \dot{\mathbf r}_I&=\mathbf v_I,
 &\dot{\mathbf v}_I&=m^{-1}\bigl(\mathcal R(q)\mathbf F_B+
                              \mathbf F_{I,\mathrm{ext}}\bigr),\notag\\
 \dot q&=\tfrac12 q\otimes(0,\boldsymbol\omega_B),
 &\dot{\boldsymbol\omega}_B
 &=J^{-1}\bigl(\boldsymbol\tau_B-
                  \boldsymbol\omega_B\times J\boldsymbol\omega_B\bigr).
 \label{gnc:simulation-rhs}
\end{align}
Thus $\dim\mathcal X=12$ although $z$ stores $13$ real components;
$q\sim-q$ gives the same physical state. The moment $\boldsymbol\tau_B$
is about the center of mass; a force at body offset $\boldsymbol\rho_B$
contributes $\boldsymbol\rho_B\times\mathbf F_B$.

For $\mathbf F_B=\mathbf F_{I,\mathrm{ext}}=\boldsymbol\tau_B=0$,
\begin{align}
 \mathbf r_I(t)&=\mathbf r_I(0)+t\mathbf v_I(0),
 &\mathbf p_I&=m\mathbf v_I=\mathrm{const},\notag\\
 \mathbf H_I&=\mathcal R(q)J\boldsymbol\omega_B=\mathrm{const},
 &E&=\tfrac12m\lVert\mathbf v_I\rVert^2+
            \tfrac12\boldsymbol\omega_B^TJ\boldsymbol\omega_B
            =\mathrm{const}.\label{gnc:simulation-invariants}
\end{align}
Here $\mathbf H_I$ is spin about the center of mass; total angular momentum
about the inertial origin is $\mathbf r_I\times\mathbf p_I+\mathbf H_I$.
Generally $\dot{\boldsymbol\omega}_B\ne0$ even with no external wrench.
For distinct principal moments, steady spin about the largest or smallest
principal moment is stable; spin about the intermediate moment is unstable
(Goldstein \S5.6, pp.~200--208 \cite{dgdt:GPS2001}).

\subsection{Midpoint momentum and a Cayley attitude step}
For step size $h>0$, let $\boldsymbol\ell_n=J\boldsymbol\omega_{B,n}$.
For a body-fixed wrench held constant over the step, solve
\begin{align}
 \bar{\boldsymbol\ell}&=\tfrac12(\boldsymbol\ell_n+
                                              \boldsymbol\ell_{n+1}),
 &\bar{\boldsymbol\omega}&=J^{-1}\bar{\boldsymbol\ell},\notag\\
 \boldsymbol\ell_{n+1}-\boldsymbol\ell_n
 &=h\bigl(\boldsymbol\tau_B-
               \bar{\boldsymbol\omega}\times\bar{\boldsymbol\ell}\bigr),
 &u&=\frac{(1,\tfrac h2\bar{\boldsymbol\omega})}
             {\sqrt{1+\tfrac{h^2}4\lVert\bar{\boldsymbol\omega}\rVert^2}},
 \notag\\
 q_{n+1}&=q_n\otimes u,
 &R_{n+1}&=R_n\,
 \underbrace{(I-\tfrac h2\widehat{\bar{\boldsymbol\omega}})^{-1}
             (I+\tfrac h2\widehat{\bar{\boldsymbol\omega}})}_
             {\operatorname{cay}(h\widehat{\bar{\boldsymbol\omega}})}.
 \label{gnc:midpoint-cayley}
\end{align}
The nonlinear solve uses its residual, rather than a fixed number of
iterations, to decide convergence. Let
$q_{1/2}=(q_n+q_{n+1})/\lVert q_n+q_{n+1}\rVert$ using these continuous
quaternion lifts. Then
\begin{align*}
 \mathbf a_{1/2}
 &=m^{-1}\bigl(\mathcal R(q_{1/2})\mathbf F_B+
                                      \mathbf F_{I,\mathrm{ext}}\bigr),\\
 \mathbf v_{n+1}&=\mathbf v_n+h\mathbf a_{1/2},\qquad
 \mathbf r_{n+1}=\mathbf r_n+\tfrac h2(\mathbf v_n+\mathbf v_{n+1}).
\end{align*}
In exact arithmetic, for zero wrench,
\begin{align*}
 T_{\mathrm{rot},n+1}-T_{\mathrm{rot},n}
 &=\bar{\boldsymbol\omega}\cdot
                 (\boldsymbol\ell_{n+1}-\boldsymbol\ell_n)=0,\\
 (I+\tfrac h2\widehat{\bar{\boldsymbol\omega}})
       \boldsymbol\ell_{n+1}
 &=(I-\tfrac h2\widehat{\bar{\boldsymbol\omega}})
       \boldsymbol\ell_n
 \ \Longrightarrow\ R_{n+1}\boldsymbol\ell_{n+1}
                         =R_n\boldsymbol\ell_n.
\end{align*}
Floating-point conservation is limited by the nonlinear residual and
roundoff. This is a second-order trajectory approximation, not the exact
flow: even spherical inertia gives the step angle
$2\arctan(h\lVert\boldsymbol\omega\rVert/2)$ rather than
$h\lVert\boldsymbol\omega\rVert$.
Check both invariant drift and attitude/position convergence against an
independent reference; invariant preservation alone cannot bound phase error.

\subsection{What the model resolves; what to read next}
Equations~\eqref{gnc:simulation-rhs} suffice for an ideal rigid body in
force-free, torque-free space. A browser clock requests elapsed time;
the backend advances fixed physical steps. Applied controls are explicit
external wrenches; releasing them returns to free coast.

For dynamics, read Goldstein \S\S5.1, 5.3--5.6 (pp.~184--208)
\cite{dgdt:GPS2001}; for an engineering connection, Wie
\S5.5 (p.~339) and \S\S6.1--6.7 (pp.~349--371)
\cite{BWie2008}. Wie \S\S7.3--7.7 add quaternion feedback, momentum
management and jet selection; Chapters~8--9 add flexibility and slosh.
Spacecraft fidelity additionally requires measured mass properties and
appropriate environment/actuator models. Numerical verification does not
by itself validate those physical assumptions. NASA's six-DOF check cases
\cite{NASA6DOF2015} provide a later independent benchmark target.


\section{Which metric makes rigid-body motion geodesic?}\label{gnc:sec:geodesic}
The submanifold proof does not by itself specify the body's dynamics.
For a rigid body with a material point $P$ fixed, assume a constant
positive-definite body inertia tensor $J_P$ about that point. The kinetic
energy defines a Riemannian metric on $SO(3)$ by
\begin{equation}\label{gnc:kinetic-metric}
 g_R(R\widehat{\mathbf a},R\widehat{\mathbf b})
 =\mathbf a^TJ_P\mathbf b,\qquad
 T=\tfrac12g_R(\dot R,\dot R)
  =\tfrac12\boldsymbol\omega_B^TJ_P\boldsymbol\omega_B.
\end{equation}
If $P$ is the center of mass, $J_P$ is the $J_B$ used above.
The metric inherited directly from the Euclidean matrix space gives
\[
 \operatorname{tr}\bigl((R\widehat{\mathbf a})^T
                         R\widehat{\mathbf b}\bigr)
 =\operatorname{tr}(\widehat{\mathbf a}^{\,T}\widehat{\mathbf b})
 =2\mathbf a\cdot\mathbf b.
\]
It is proportional to the kinetic metric only for spherical inertia
$J_P=jI$. The embedding in $\mathbb R^9$ therefore does not determine
the metric of a general rigid body.

In any local coordinates $q^a$ on $SO(3)$, write
$T=\tfrac12g_{ab}(q)\dot q^a\dot q^b$. For a smooth time-independent
potential $V(q)$, the Lagrangian $L=T-V$ has
\[
 \frac{\partial L}{\partial\dot q^a}=g_{ab}\dot q^b,\qquad
 \frac{\partial L}{\partial q^a}
 =\tfrac12\frac{\partial g_{bc}}{\partial q^a}\dot q^b\dot q^c
  -\frac{\partial V}{\partial q^a}.
\]
Substituting into the Euler--Lagrange equations gives
\begin{equation}\label{gnc:metric-euler-lagrange}
 \frac{d}{dt}(g_{ab}\dot q^b)
 -\tfrac12\frac{\partial g_{bc}}{\partial q^a}\dot q^b\dot q^c
 =-\frac{\partial V}{\partial q^a}.
\end{equation}
For constant $V$, Eq.~\eqref{gnc:metric-euler-lagrange} is the
affinely parameterized geodesic equation of $g$ in physical time.
For a general $V$, the right-hand side is a force term.

At fixed energy $E=T+V$, on the region $E>V$, Jacobi's principle
describes the same configuration paths as geodesics of
\begin{equation}\label{gnc:jacobi-metric}
 g^{(E)}=2(E-V)g.
\end{equation}
Indeed, $ds_g/dt=\sqrt{g_R(\dot R,\dot R)}=\sqrt{2(E-V)}$, so
\begin{equation}\label{gnc:jacobi-time}
 ds_E=\sqrt{2(E-V)}\,ds_g,\qquad
 \frac{ds_E}{dt}=2(E-V).
\end{equation}
The reduced action $\int p_a\,dq^a$ equals $\int ds_E$ along
energy-$E$ parametrizations, because
$p_a\,dq^a=g_{ab}\dot q^b\dot q^a\,dt=2T\,dt=ds_E$.
Its stationary paths are therefore geodesics for $g^{(E)}$.
Frankel \cite{TFrankel2012}, Section~10.2c, Theorem~(10.19),
uses $d\rho=\sqrt{E-V}\,ds_g$, so his metric is
$(E-V)g=\tfrac12g^{(E)}$; the constant factor changes arc length
but not the geodesic paths. The metric degenerates at $E=V$, and
physical time is generally not its affine geodesic parameter.

If $E>\max_{SO(3)}V$, the Jacobi metric is positive everywhere on
the compact configuration manifold. Closed-geodesic results can then
be used to establish periodic motions, as in Frankel's Section~10.2d.
A closed geodesic returns with the same tangent direction; returning
to the same configuration alone is insufficient.
Its physical period is recovered by integrating
$dt=ds_E/[2(E-V)]$ around the closed geodesic.

\section{Reading geometric mechanics}
Holm, Schmah, and Stoica \cite{HSS2009} provide the geometric foundation.
Start with Chapter~1 (pp.~3--42), then manifolds in Chapter~2
(pp.~43--98) and geometry in Chapter~3 (pp.~99--154). Read mechanics on
manifolds in Chapter~4 (pp.~155--186), Lie groups and algebras in
Chapter~5 (pp.~187--208), and group actions in Chapter~6 (pp.~209--240).
The first destination is Sections~7.1--7.3 (pp.~241--260): rigid-body
mechanics and Euler--Poincar\'e reduction. Sections~7.4--7.5
(pp.~261--280) introduce the heavy top and advected parameters.
Chapters~8--9 (pp.~281--324) develop momentum maps and Lie--Poisson
reduction. Printed page numbers refer to the 2009 Oxford edition.
The page-reconciled transcription is \texttt{resolved.md} in that book's
\texttt{ocr-compare/reconciled} directory. These notes do not copy the book.

Derive Eq.~\eqref{gnc:euler} again from the reduced variational principle
when studying Chapter~7. Compare that derivation with the momentum balance
above. Keep frame conventions explicit throughout. Guidance, estimation,
feedback control, aerodynamic models, and open-system propulsion balances
require further material beyond these geometric-mechanics chapters.
